\section{预备知识：概率论与数理统计基础速览}

本章是后面所有内容的"地基"。统计计算的核心思想只有一句话：
\textbf{用计算机生成大量随机样本，再用样本的频率/均值去近似概率/期望。}
要看懂这件事，需要先把下面几个基础概念弄扎实。

\subsection{随机变量、分布函数与密度函数}

\begin{kbox}{核心概念}
\begin{itemize}[leftmargin=2em]
  \item \textbf{分布函数（CDF）}：$F(x)=\Pr(X\le x)$。单调不减、右连续，$F(-\infty)=0$，$F(+\infty)=1$。任何随机变量都有分布函数。
  \item \textbf{密度函数（pdf）}：连续型随机变量满足 $F(x)=\int_{-\infty}^{x}f(t)\d t$，即 $f(x)=F'(x)$。直观理解：$f(x)\Delta x \approx \Pr(x< X\le x+\Delta x)$，密度大的地方取值"更密集"。
  \item \textbf{分布律（pmf）}：离散型随机变量用 $p_k=\Pr(X=x_k)$ 描述，$\sum_k p_k=1$。
\end{itemize}
\end{kbox}

\subsection{期望、方差与协方差}

\[
\E(X)=\int xf(x)\d x \quad(\text{连续型}),\qquad
\E(X)=\sum_k x_k p_k \quad(\text{离散型}).
\]
\[
\Var(X)=\E\big[(X-\E X)^2\big]=\E(X^2)-(\E X)^2 .
\]

常用性质（考试与推导中反复用到）：
\begin{itemize}[leftmargin=2em]
  \item 线性性：$\E(aX+bY)=a\E X+b\E Y$（\textbf{无需独立}）。
  \item $\Var(aX+b)=a^2\Var(X)$。
  \item $\Var(X+Y)=\Var X+\Var Y+2\Cov(X,Y)$；若 $X,Y$ 独立，则 $\Cov(X,Y)=0$。
  \item $\Cov(X,Y)=\E(XY)-\E X\,\E Y$；相关系数 $\rho=\Cov(X,Y)/\sqrt{\Var X\,\Var Y}\in[-1,1]$。
\end{itemize}

\begin{rmk}
第5章方差缩减技术的全部原理就建立在
$\Var\!\big(\frac{X+Y}{2}\big)=\frac14(\Var X+\Var Y+2\Cov(X,Y))$
这个公式上：若能让 $\Cov(X,Y)<0$，平均后的方差就比独立样本更小。
\end{rmk}

\subsection{常用分布速查表}

\begin{center}
\renewcommand{\arraystretch}{1.5}
\begin{tabular}{|l|l|l|l|}
\hline
分布 & 密度/分布律 & 期望 & 方差\\\hline
均匀 $U(a,b)$ & $f=\frac{1}{b-a},\ a<x<b$ & $\frac{a+b}{2}$ & $\frac{(b-a)^2}{12}$\\\hline
指数 $\mathrm{Exp}(\lambda)$ & $f=\lambda e^{-\lambda x},\ x>0$ & $1/\lambda$ & $1/\lambda^2$\\\hline
正态 $N(\mu,\sigma^2)$ & $f=\frac{1}{\sqrt{2\pi}\sigma}e^{-\frac{(x-\mu)^2}{2\sigma^2}}$ & $\mu$ & $\sigma^2$\\\hline
Gamma$(\alpha,\gamma)$ & $f=\frac{\gamma^{\alpha}}{\Gamma(\alpha)}x^{\alpha-1}e^{-\gamma x},\ x>0$ & $\alpha/\gamma$ & $\alpha/\gamma^2$\\\hline
Beta$(\alpha,\beta)$ & $f=\frac{\Gamma(\alpha+\beta)}{\Gamma(\alpha)\Gamma(\beta)}x^{\alpha-1}(1-x)^{\beta-1},\ 0<x<1$ & $\frac{\alpha}{\alpha+\beta}$ & $\frac{\alpha\beta}{(\alpha+\beta)^2(\alpha+\beta+1)}$\\\hline
二项 $B(n,p)$ & $\Pr(X=k)=\binom{n}{k}p^k(1-p)^{n-k}$ & $np$ & $np(1-p)$\\\hline
泊松 $P(\lambda)$ & $\Pr(X=k)=\frac{\lambda^k}{k!}e^{-\lambda}$ & $\lambda$ & $\lambda$\\\hline
几何 $Geo(p)$ & $\Pr(X=k)=(1-p)^{k-1}p$ & $1/p$ & $\frac{1-p}{p^2}$\\\hline
\end{tabular}
\end{center}

\textbf{几个重要关系}（生成随机数时经常用来"借力"）：
\begin{itemize}[leftmargin=2em]
  \item $\mathrm{Exp}(\lambda)=\mathrm{Gamma}(1,\lambda)$；Gamma 可加性：独立的 $\mathrm{Gamma}(\alpha_i,\gamma)$ 之和服从 $\mathrm{Gamma}(\sum\alpha_i,\gamma)$。
  \item $\chi^2(n)=\mathrm{Gamma}(n/2,1/2)$；若 $Z_i\overset{iid}{\sim}N(0,1)$，则 $\sum_{i=1}^n Z_i^2\sim\chi^2(n)$。
  \item 若 $\xi\sim\mathrm{Gamma}(\alpha,\gamma)$，则 $\xi/a\sim\mathrm{Gamma}(\alpha,a\gamma)$（尺度变换）。
  \item 若 $Z\sim N(0,1)$，则 $\mu+\sigma Z\sim N(\mu,\sigma^2)$。
\end{itemize}

\subsection{三大抽样分布与统计量}

设 $X_1,\dots,X_n\overset{iid}{\sim}N(\mu,\sigma^2)$，
样本均值 $\bar X=\frac1n\sum X_i$，样本方差 $S^2=\frac{1}{n-1}\sum(X_i-\bar X)^2$。则
\[
\bar X\sim N\!\Big(\mu,\frac{\sigma^2}{n}\Big),\qquad
\frac{(n-1)S^2}{\sigma^2}\sim\chi^2(n-1),\qquad
\frac{\bar X-\mu}{S/\sqrt n}\sim t(n-1),
\]
且 $\bar X$ 与 $S^2$ 独立。两个独立卡方变量标准化之比服从 $F$ 分布：
$\dfrac{\chi^2(m)/m}{\chi^2(n)/n}\sim F(m,n)$。

\subsection{大数定律与中心极限定理（蒙特卡洛方法的理论基础）}

\begin{kbox}{为什么"用样本均值估期望"是合理的？}
设 $X_1,X_2,\dots$ 独立同分布，$\E X_i=\mu$，$\Var X_i=\sigma^2<\infty$。
\begin{itemize}[leftmargin=2em]
  \item \textbf{大数定律（LLN）}：$\bar X_n=\frac1n\sum_{i=1}^n X_i\xrightarrow{P}\mu$。\\
  即样本量足够大时，样本均值几乎必然接近真实期望——这保证了蒙特卡洛估计的\textbf{相合性}。
  \item \textbf{中心极限定理（CLT}）：$\dfrac{\bar X_n-\mu}{\sigma/\sqrt n}\xrightarrow{d}N(0,1)$。\\
  即误差约为 $O(1/\sqrt n)$ 量级，且近似服从正态分布——这给出了蒙特卡洛估计的\textbf{精度}与置信区间。
\end{itemize}
\textbf{重要结论}：蒙特卡洛误差的标准差为 $\sigma/\sqrt n$。想把精度提高 10 倍，样本量需要扩大 100 倍；或者保持 $n$ 不变、想办法\textbf{减小 $\sigma$}——这正是第 5 章方差缩减的动机。
\end{kbox}

\subsection{点估计、置信区间与假设检验回顾}

\begin{itemize}[leftmargin=2em]
  \item \textbf{点估计}：用统计量 $\hat\theta(X_1,\dots,X_n)$ 估计未知参数 $\theta$。构造方法主要有\textbf{矩估计法}和\textbf{极大似然估计法（MLE）}。评价标准：
  \begin{itemize}
    \item 无偏性：$\E(\hat\theta)=\theta$；偏差 $\mathrm{Bias}(\hat\theta)=\E(\hat\theta)-\theta$。
    \item 有效性：同为无偏估计时方差更小者更有效。
    \item 相合性：$n\to\infty$ 时 $\hat\theta\xrightarrow{P}\theta$。
    \item 综合指标——均方误差：$\mathrm{MSE}(\hat\theta)=\E(\hat\theta-\theta)^2=\mathrm{Bias}^2+\Var(\hat\theta)$。
  \end{itemize}
  \item \textbf{置信区间}：随机区间 $[\hat\theta_L,\hat\theta_U]$ 满足 $\Pr(\hat\theta_L\le\theta\le\hat\theta_U)=1-\alpha$。$1-\alpha$ 称置信水平。构造的通用套路是找\textbf{枢轴量}（含 $\theta$ 但分布已知的量）。
  \item \textbf{假设检验}：原假设 $H_0$ vs 备择假设 $H_1$。
  \begin{itemize}
    \item \textbf{第一类错误}（弃真）：$H_0$ 为真却拒绝，概率记 $\alpha$（显著性水平）。
    \item \textbf{第二类错误}（取伪）：$H_0$ 为假却接受，概率记 $\beta$。
    \item \textbf{功效（power）}：$1-\beta=\Pr(\text{拒绝 }H_0\mid H_1\text{ 成立})$，衡量检验"识别假原假设"的能力。
  \end{itemize}
  \item \textbf{似然函数}：样本 $x_1,\dots,x_n$ 给定后，$L(\theta)=\prod_{i=1}^n f(x_i;\theta)$。MLE 就是使 $L(\theta)$（等价地 $\log L$）最大的 $\theta$。第 7 章 EM 算法处理的是"数据不完整时没法直接最大化似然"的情形。
\end{itemize}
